\section{RT-1 的 Scaling 与 Generalization 证据}

训练任务上的平均 success rate 不能说明机器人能否进入新厨房、忽略 distractor 或组合未见指令。RT-1 因此沿 seen tasks、unseen tasks、distractors 和 backgrounds 四个轴评估，并进一步测试多个变化因素同时出现时的性能。

\subsection{Visual 与 Semantic Diversity}

上一章只证明 RT-1 能在实时约束下执行多任务控制，还没有回答这些能力是否依赖训练厨房、固定物体与模板指令。本节先把 generalization 拆成四种可区分的分布变化；读图时应分别比较 seen、unseen、distractor 与 background 条件，而不是只看最高柱或总体平均。这里的结果能支持相对 baseline 的优势，却不能证明模型获得了开放世界语义理解。

slide 17 左侧四组柱子比较 RT-1、GATO、BC-Z 与 BC-Z XL，右侧展示背景、物体和干扰变化。RT-1 在四组条件下整体更强，但 unseen / background 条件仍明显低于 seen tasks，说明“最好 baseline”不等于已经解决开放世界泛化。

\lecturefigure{slide-17-rt1-visual-semantic-diversity.jpg}{RT-1 在 seen、unseen、distractor 与 background 变化上的成功率}{官方录像恢复幻灯片 V017；Brohan et al., 2022}

\begin{knowledgebox}{读图：四根柱子测的是不同失败模式}
Seen tasks 主要检验训练分布拟合；unseen tasks 检验语言与行为组合；distractors 检验注意力是否被无关物体带偏；backgrounds 检验视觉分布变化。四项不能压成一个平均数，因为不同 architecture component 对不同轴的贡献不同。
\end{knowledgebox}

\subsection{多因素变化比单因素更难}

前一张图把变化因素逐项隔离，便于定位模型在哪一类 shift 上退化；真实部署却很少一次只改变一个变量。本节因此转向 compound shift，问题是背景、布局、物体位置、视角和语言同时变化时，单因素优势能否继续保持。看图时先比较 level 1 到 level 3 的递进下降，再观察 unseen kitchen 的具体差异；这组演示说明联合变化更难，但不能把少量成功案例当作任意新环境都可迁移的证据。

slide 18 把柱状结果与 unseen kitchen 视频并排：机器人不仅换背景，还可能同时遇到新物体位置、不同桌面、语言变化和视角差异。stacked bars 的 level 1/2/3 表示越来越多因素同时变化，RT-1 虽领先，但最高难度成功率仍有限。

\lecturefigure{slide-18-rt1-variation-unseen-kitchen.jpg}{多因素 generalization：在 unseen kitchen 中同时改变背景、布局和任务}{官方录像恢复幻灯片 V018；Stanford CS25 Lecture 15}

第二个 demo 使用自然语言变体，例如不再重复训练模板，而以更自由的描述要求机器人行动。它展示 language conditioning 的价值，也暴露任务语义仍依赖已有 skill coverage：模型可以理解换一种说法，却不能凭语言创造训练数据从未包含的物理能力。

\lecturefigure{slide-19-rt1-language-variation-demo.jpg}{语言变化演示：在多因素视觉变化中执行不同措辞的任务}{官方录像恢复幻灯片 V019；Stanford CS25 Lecture 15}

\teachervoice{讲者把这些视频当作定量图的补充，而不是“挑最好案例”。他反复提醒多因素 generalization 远难于单变量测试；机器人在新厨房工作几次，并不等于已经学会任意家庭环境。}

\subsection{跨 Distribution 的数据混合}

RT-1 还混入 simulation / sim-to-real、bin-picking 和不同 embodiment 数据。图中左侧是数据来源，右侧柱状条报告它们对 targeted evaluations 的相对改善。这里的结论是高容量模型能从异质数据中提取共享信息，而不是所有数据无条件正迁移。

\lecturefigure{slide-20-rt1-diverse-data-distributions.jpg}{训练在多种 data distributions：everyday robots、simulation 与 bin picking}{官方录像恢复幻灯片 V020；Stanford CS25 Lecture 15}

\begin{warningbox}{Data interoperability 有边界}
不同机器人拥有不同相机、动作空间和动力学。若没有 embodiment identifier、action normalization 或共享语义，简单拼接数据会产生冲突。RT-1 的结果是“部分异质数据可以帮助特定分布”，不是“任意 robot dataset 都能直接合并”。
\end{warningbox}

\subsection{Task Diversity 比 Raw Size 更关键}

data scaling slide 的横轴是保留的数据比例，纵轴是 success rate；蓝点是 seen tasks，黑点是 generalization。绿色箭头表示减少总数据量，紫色箭头表示在保留较多 episode 的同时减少 task diversity。后者导致的下降更陡，说明重复同一任务不能替代覆盖更多任务。

\lecturefigure{slide-21-rt1-data-scaling.jpg}{Scaling with data：减少 task diversity 比等量减少 episode 更伤泛化}{官方录像恢复幻灯片 V021；Brohan et al., 2022}

\begin{importantbox}{机器人数据的四个“规模”}
\begin{enumerate}
  \item \textbf{Episode count}：轨迹总数；
  \item \textbf{Task diversity}：不同技能与目标组合；
  \item \textbf{Environment diversity}：背景、布局、物体与 lighting；
  \item \textbf{Embodiment diversity}：不同机器人和 action space。
\end{enumerate}
“数据更大”只有说明是哪一种规模，才具有工程意义。
\end{importantbox}

\subsection{Ablation：不是单一 Transformer 效应}

ablation 图以完整 RT-1 为零点，向下表示成功率损失。GATO、BC-Z、BC-Z XL、去掉大模型、去掉 pretraining、改连续/自回归 action、去掉 history 或 Transformer，分别在四种评价轴产生不同下降。结果说明 architecture、pretraining、history 和 action representation 都参与性能，而非只靠参数量。

\lecturefigure{slide-22-rt1-design-ablations.jpg}{RT-1 design ablations：各组件对四类 generalization 的相对影响}{官方录像恢复幻灯片 V022；Brohan et al., 2022}

\begin{knowledgebox}{读图：不要只找最长的柱}
每种颜色对应不同 evaluation axis，某组件可能主要帮助 unseen tasks，却对 seen tasks 影响较小。`w/o transformer` 或 `w/o pretraining` 的损失不能被单独解释为 causal law，因为替代结构、训练稳定性与容量可能同时改变。Ablation 的价值是定位脆弱轴，而不是给组件排绝对名次。
\end{knowledgebox}

\subsection{本章小结}

RT-1 的证据显示 Transformer-style token interface 能在实时预算下提升多任务控制与多个泛化轴，但能力仍随场景复杂度下降。最强数据结论是 task diversity 高于简单 episode repetition；最强系统结论是多组件共同作用，不能把成功归因于“用了 Transformer”一句话。

